\chapter{Introducción}
\label{ch:introduccion}

\section{Planteamiento del Problema}

En contraste con las funciones globales, el uso de B-splines (polinomios a trozos finitos) introduce un marco numérico de alta precisión libre de dependencias lineales que es inherentemente agnóstico a la física subyacente. Al ser un constructo puramente matemático, comúnmente empleado en los métodos de elementos finitos, proporciona una flexibilidad representativa extraordinaria. Esta base finita garantiza un soporte riguroso capaz de aproximar orbitales con precisión casi ilimitada (autores como \cite{Bachau2001, deBoor1978} aclaman obtener precisión de máquina), al mismo tiempo que discretiza eficazmente el continuo atómico, permitiendo estudios profundos sobre fenómenos de fotoionización y correlación electrónica \cite{Johnson2007}. De hecho, la robustez de los B-splines no se limita a la ecuación radial de Schrödinger, sino que ha demostrado ser estable en la resolución de la ecuación relativista de Dirac-Hartree-Fock para elementos pesados \cite{Johnson1988}.

El objetivo principal de esta tesis es desarrollar, implementar e iterar un solucionador numérico \textit{ab-initio} fundamentado en B-splines para sistemas atómicos multielectrónicos. En particular, nos enfocamos en el método Restricted Open-Shell Hartree-Fock (ROHF) para estudiar con extrema fidelidad la secuencia isoelectrónica $np^2$, conformada por el Carbono ($2p^2$), Silicio ($3p^2$), Germanio ($4p^2$) y Estaño ($5p^2$). 

A medida que descendemos por el grupo 14 de la tabla periódica, el aumento en la carga nuclear $Z$ y el efecto de apantallamiento de las capas internas cerradas provocan un rápido incremento en las fuerzas relativistas efectivas experimentadas por los electrones de valencia. El reto es demostrar que nuestra base de B-splines no solo obtiene energías autoconsistentes exactas a nivel no-relativista, sino que permite calcular con precisión el parámetro de acoplamiento espín-órbita radial ($\zeta_{np}$) y evaluar la matriz del Hamiltoniano de Breit-Pauli. Este análisis culminará demostrando el desglose de la aproximación tradicional del acoplamiento Russell-Saunders ($LS$) en el Germanio y el Estaño, revelando una transición contundente hacia un régimen de acoplamiento intermedio.

Es fundamental enmarcar que, dada la naturaleza de este trabajo como un proyecto de tesis de licenciatura, el objetivo central no radica en postular un desarrollo teórico original para el estado del arte de la física atómica. Por el contrario, la presente investigación posee un carácter netamente formativo y exploratorio. La implementación desde cero de este motor \textit{ab-initio} busca demostrar un dominio analítico y numérico profundo de la mecánica cuántica multielectrónica y la fenomenología relativista. Al aplicar de manera formal la teoría espectroscópica de \cite{Cowan1981}, se establece una base metodológica que avala la capacidad técnica y conceptual requerida para abordar problemas avanzados en futuros estudios de posgrado.

\section{Estrategia de Modelado: De Hartree-Fock al Acoplamiento Intermedio}

Para comprender la metodología empleada y los resultados obtenidos en este tratado, es fundamental asimilar el flujo lógico del modelo físico, el cual evoluciona desde una perspectiva teórica básica hacia una representación empírica cada vez más sofisticada.

El desarrollo se cimenta en primer lugar sobre el esquema tradicional uniconfiguracional de Hartree-Fock. En esta etapa, el problema de muchos cuerpos se aborda mediante la aproximación de campo medio, resolviendo iterativamente las ecuaciones integro-diferenciales para alcanzar la convergencia del Campo Autoconsistente (SCF). Los orbitales obtenidos en este límite, aunque útiles para una descripción estática, adolecen de precisión empírica dado que ignoran efectos de correlación y relativistas cruciales para los átomos pesados.

Una vez establecidos y convergidos estos resultados del SCF, la estrategia de modelado consiste en utilizar dichos orbitales optimizados como una base confiable para introducir sistemáticamente correcciones físicas. De este modo, procedemos a mejorar nuestro modelo incorporando los siguientes componentes fenomenológicos:
\begin{itemize}
    \item \textbf{Interacción de Configuraciones (CI):} Para recuperar la energía de correlación estática omitida en el determinante de Slater inicial.
    \item \textbf{Acoplamiento Espín-Órbita (SO):} Integrando el Hamiltoniano de Breit-Pauli sobre nuestros orbitales para capturar el desdoblamiento magnético característico del régimen de acoplamiento intermedio.
    \item \textbf{Potencial de Polarización ($V_{\text{pol}}$):} Una corrección semiempírica que modela la interacción del core inerte con los electrones de valencia, vital para afinar las predicciones energéticas en elementos pesados como el Germanio y el Estaño.
\end{itemize}

A través de esta escalada metodológica, transitamos del modelo estricto de partículas independientes hacia un modelo robusto que reproduce con alta fidelidad las observaciones espectroscópicas empíricas.
\section{Estructura del Documento}

Este trabajo se ha estructurado de la siguiente manera:
En el \textbf{Capítulo 2}, se presenta el marco teórico subyacente, introduciendo la mecánica del determinante de Slater, las integrales de repulsión de Coulomb y el ensamble del Hamiltoniano de muchos cuerpos.
El \textbf{Capítulo 3} describe el método numérico, detallando la fundamentación teórica de los B-splines y la formulación matricial de las ecuaciones de Hartree-Fock (método de Galerkin).
El \textbf{Capítulo 4} expone la validación numérica de la implementación (convergencia virial, colapso de potenciales y energías de correlación espín-órbita). 
El \textbf{Capítulo 5} discute el contraste físico a lo largo de la secuencia $np^2$ y la transición al acoplamiento intermedio.
Por último, el \textbf{Capítulo 6} sintetiza las conclusiones principales derivadas de esta investigación.

Para facilitar la legibilidad del texto principal y mantener un hilo conductor enfocado en la física del problema, los desarrollos formales más extensos y técnicos se han delegado a una serie de apéndices. El \textbf{Apéndice A} precisa el sistema de unidades atómicas; el \textbf{Apéndice B} desarrolla la teoría general de momento angular; el \textbf{Apéndice C} expone el formalismo avanzado de operadores tensoriales irreducibles; el \textbf{Apéndice D} presenta los fundamentos algorítmicos de la base de B-splines; el \textbf{Apéndice E} presenta la formulación de Galerkin y la resolución de la ecuación de Poisson radial que sustenta el cálculo de los potenciales de interacción electrónica; y, finalmente, el \textbf{Apéndice F} documenta los métodos numéricos auxiliares empleados durante la implementación.

%%% Local Variables:
%%% mode: LaTeX
%%% TeX-master: "../main"
%%% End:
